% Source: https://artofproblemsolving.com/community/c3h1693928_geometry
% Stable identity: aops-topic-1693928-q01
Let $R$ and $r$ be the radii of $Y$ and $X$, and $Q$ the centre of $X$. Since the circles touch internally at $P$, the points $O$, $Q$, $P$ are collinear and $OQ=R-r$.

\textbf{The ratio of the radii.} The homothety with centre $P$ and ratio $\dfrac {R}{r}$ maps circle $X$ onto circle $Y$ (it maps $Q$ to $O$, since $PO=R=\dfrac {R}{r}\cdot PQ$). It maps the point $Z$ of $X$ to the point of $Y$ on ray $PZ$ other than $P$, which is $M$. Hence $PM=\dfrac {R}{r}\,PZ$. But $PM=PZ+ZM=3PZ$, so $R=3r$ and $OQ=2r$.

\textbf{The angle at $O$.} Let $T$ be the point where $MN$ touches $X$; then $QT\perp MN$ and $QT=r$. In right triangle $OTQ$,
\[
\sin\angle QOT=\dfrac{QT}{OQ}=\dfrac{r}{2r}=\dfrac{1}{2},\qquad \angle QOT=30^\circ .
\]
So the ray $OP$ (which contains $Q$) makes an angle of $30^\circ$ with the line $MN$: one of $\angle PON$, $\angle POM$ is $30^\circ$ and the other $150^\circ$. Since $PM>PN$ and a longer chord subtends a larger central angle, $\angle PON=30^\circ$.

\textbf{The answer.} The inscribed angle $\angle PMN$ stands on the arc $PN$, so it is half the central angle:
\[
\angle PMN=\dfrac{1}{2}\angle PON=\boxed{15^\circ}.
\]

\begin{asy}
usepackage("amsmath");
size(7cm);
defaultpen(fontsize(10pt));
real R=3,r=1;
pair O=(0,0),P=R*dir(30),Q=(R-r)*dir(30),M=(-R,0),Nn=(R,0);
pair Z=P+(M-P)/3,T=(Q.x,0);
draw(circle(O,R),linewidth(0.9));draw(circle(Q,r),linewidth(0.9)+deepgreen);
draw(M--Nn,linewidth(0.8));draw(P--M,linewidth(0.8));
draw(O--P,blue+dashed);draw(Q--T,blue+dashed);
real q=0.18;draw(T+(q,0)--T+(q,q)--T+(0,q),blue);
draw(arc(O,0.7,0,30),red);label("$30^\circ$",O+1.05*dir(15),red);
draw(arc(M,1.3,0,15),red);label("$15^\circ$",M+1.75*dir(7.5),red);
dot(O);dot(P);dot(Q);dot(M);dot(Nn);dot(Z);dot(T);
label("$O$",O,S);label("$P$",P,NE);label("$Q$",Q,dir(-40));label("$M$",M,W);label("$N$",Nn,E);
label("$Z$",Z,N);label("$T$",T,S);
label("$Y$",R*dir(140),dir(140));label("$X$",Q+r*dir(-30),dir(-30),deepgreen);
label("\mbox{}",(min(currentpicture,user=true)+max(currentpicture,user=true))/2,black); // keeps the text after the figure black when LaTeX typesets the labels inline
\end{asy}
